\documentclass[10pt,a4paper]{amsart}
\usepackage[margin=19mm]{geometry}
\usepackage{amsmath,amssymb,mathtools}
\usepackage[scr=rsfs,cal=euler]{mathalfa}
\usepackage{xfrac}
\usepackage[unicode,hidelinks,bookmarks=false]{hyperref}
\newcommand{\ie}{i.e.\ }
\newcommand{\cf}{cf.\ }
\newcommand{\resp}{respectively\ }
\newcommand{\Subsection}{\subsection}
\providecommand{\nobreakdash}{\nobreak\hbox{-}}
\setlength{\emergencystretch}{2em}
\multlinegap0pt
\newcommand\Coefd{\frac{2}{3}}
\usepackage{enumerate}
\usepackage{amssymb}
\usepackage{tikz}
\usepackage{bm}
\newtheorem{claim}{Claim}
\newcommand\Ga{\Gamma}
\newcommand\al{\alpha}
\newcommand\be{\beta}
\newcommand\de{\delta}
\newcommand{\defeq}{\mathrel{\vcenter{\baselineskip0.5ex \lineskiplimit0pt
                     \hbox{\scriptsize.}\hbox{\scriptsize.}}}%
                     =}
\newcommand\eps{\varepsilon}
\newcommand\ga{\gamma}
\title{Boosted second moment method\\in random regular graphs}
\author{Bal\'azs Gerencs\'er, Viktor Harangi}
\date{Source: arXiv:2510.12600v3 (CC BY 4.0)}
\begin{document}
\maketitle

\noindent\emph{Source and licence.} Boosted second moment method\\in random regular graphs. Version arXiv:2510.12600v3 (CC BY 4.0), distributed by arXiv under the Creative Commons Attribution 4.0 International licence, http://creativecommons.org/licenses/by/4.0/, as recorded in the arXiv licence metadata for this exact version and shown on the abstract page https://arxiv.org/abs/2510.12600v3. Full licence notice: \texttt{licenses/arxiv-2510.12600-CC-BY-4.0.txt}.\par\smallskip

\noindent\textit{Editorial scope.} Counted unit: the article's claim cl:f\_diff (source0.tex lines 1602--1612). The article's complete original proof of it is delivered with it (source0.tex lines 1822--1972, one \texttt{proof} environment of 151 lines, which the article places away from the statement). No mathematics is written, repaired or re-derived: the statement and the proof are the article's own lines, in the article's own notation, and no retained line is substituted or re-flowed.  Retained uncounted as the source-local context the proof needs: lines 806--813; lines 1157--1161; lines 1169--1174.  Nothing was omitted from any retained span, so the package carries no omission record.  No retained line contains a citation, so the wrapper carries no bibliography and no bibliography entry.  No retained line contains a figure environment, an image-inclusion command or drawing code, so no image is delivered and the package declares no asset dependency.  Editorial formatting only, in addition: an AMS \texttt{amsart} preamble that restates the article's own declarations and macros used by the retained lines, namely the environments \texttt{claim} on separate counters, together with the standard packages those lines need.  The retained environments print in this document as Claim 1 in order of appearance, whereas the article prints the same items with the numbering of its own full document.  The complete native source of the article as distributed in the arXiv e-print for this exact version is bundled unchanged as \texttt{sources/187054/source0.tex}.
\begin{align}
\Ga(\be) 
&\defeq -(1/2-\al)+\sqrt{(1/2-\al)^2+(\al-\be)^2} 
\label{eq:Ga_first} \\
&= \frac{(\al-\be)^2}{(1/2-\al) 
+\sqrt{(1/2-\al)^2+(\al-\be)^2}} .
\label{eq:Ga_second}
\end{align}  

\begin{equation} \label{eq:h_diff2_simple}
h(1-x-\eps) - h(1-x) 
= \eps - x\eps -\frac12 \eps^2
- \sum_{k=3}^\infty \frac{(x+\eps)^k-x^k}{k(k-1)} .
\end{equation}

\begin{equation} \label{eq:h_diff1}
h(x-\eps) - h(x) 
= - \int_{x-\eps}^x h'(z) \,\mathrm{d}z 
\leq -\eps h'(x) = \eps(1+\log x) 
\quad \text{for any } 0 \leq \eps \leq x \leq 1.
\end{equation}

\begin{claim} \label{cl:f_diff}
Let $d \geq 318$ and $\al \leq (2\log d)/d$. Suppose that $\be=\al-\de$ for some $0\leq \de \leq \frac{3}{2d}$. Then 
%
\begin{equation*}
f_{d,\al}(\be) - \underbrace{f_{d,\al}(\al)}_{=\varphi_d(\al)} 
\leq 2 h(\de) 
+ \big( -d\al + \log(\al) + 2 \big) \de 
+ \Coefd d \de^2 .
\end{equation*}
%
\end{claim}

\begin{proof}[Proof of Claim~\ref{cl:f_diff}]
For the sake of neater formulas, we will simply \underline{write $\ga$ for $\Ga(\be)$} throughout the proof, where $\be=\al-\de$ as in the statement. With this mind, we can write $f_{d,\al}(\be)$ as follows:
%
\begin{align*}
f_{d,\al}(\be) 
&= \frac{d}{2} \bigg( 2h(\be)+2h(\ga)
+4 h(\al-\be-\ga) + h(1-4\al+2\be+2\ga) \bigg) \\
&- (d-1) \bigg( h(\be)+ 2h(\al-\be) 
+ h(1-2\al+\be) \bigg) \\
&= h(\al-\de) + d h(\ga) 
+ 2d \bigg( h(\de-\ga) - h(\de) \bigg) + 2h(\de) \\
&+ \frac{d}{2} h\big(1-2\al-2(\de-\ga)\big) 
- (d-1) h(1-\al-\de) ,
\end{align*}  
%
while \eqref{eq:Ga_second} translates to 
%
\begin{equation} \label{eq:ga_bounds}
\ga= \frac{\de^2}{\sqrt{\big(\frac12-\al\big)^2+\de^2}+\big(\frac12-\al\big)} .
\end{equation}
%
Note that we have the following bounds for $\ga$:
%
\begin{equation} \label{eq:ga_bnd}
\de^2 
\leq \frac{1}{1-\al} \de^2
\leq \ga 
\leq \frac{1}{1-2\al} \de^2 .
\end{equation}
%

Now we turn our attention to the difference in question: fix $d,\al$ and let 
%
\begin{equation} \label{eq:B_de}
B_\de \defeq f_{d,\al}(\be)-f_{d,\al}(\al) . 
\end{equation}
%
For brevity we use the notation $\de' \defeq \de-\ga$. Then 
%
\begin{align*}
B_\de 
&= \bigg( h(\al-\de) - h(\al) \bigg) + d h(\ga) 
+ 2d \bigg( h(\de') - h(\de) \bigg) + 2h(\de) \\
&+ \frac{d}{2} \bigg( h(1-2\al-2\de') - h(1-2\al) \bigg)
- (d-1) \bigg( h(1-\al-\de) - h(1-\al) \bigg).
\end{align*}  
%
Using \eqref{eq:h_diff1} with $x=\al$ and $\eps=\de$ we get 
%
\begin{equation} \label{eq:bnd1}
h(\al-\de) - h(\al) \leq -\de h'(\al) 
= \de \big( 1+\log \al \big) .
\end{equation}
%
Using \eqref{eq:h_diff1} again, this time with $x=\de$ and $\eps=\ga$ so that $x-\eps=\de'$, we get 
%
\[ h(\de') - h(\de) \leq -\ga h'(\de) 
= \ga \big( 1+ \log \de \big) ,\]
%
and hence 
%
\begin{equation} \label{eq:bnd2}
2d\bigg( h(\de') - h(\de) \bigg) + d h(\ga) 
\leq d \ga \big(2 + 2 \log \de  - \log \ga \big) 
\leq 2d \ga ,
\end{equation}
%
where we used $2 \log \de \leq \log \ga$ as $\de^2 \leq \ga$ according to \eqref{eq:ga_bnd}. 


Recall that we work under the assumptions $\al\leq 2(\log d)/d$ and $\de \leq \frac{3}{2d}$. Then, according to \eqref{eq:ga_bnd}, we have 
%
\[ \frac{\ga}{\de} 
\leq \frac{\de}{1-2\al}
\leq \frac{3/(2d)}{1-4(\log d)/d} 
= \frac{3/2}{d-4(\log d)} < 0.37 
\text{ for any } d \geq 15. \]
%
It follows that 
\begin{equation} \label{eq:deperde}
\frac{\de'}{\de} = 1 - \frac{\ga}{\de} 
> 1- 0.37 
> \frac{1}{\sqrt[3]{4}} .
\end{equation}
%
Using \eqref{eq:h_diff2_simple} with $x=\al$ and $\eps=\de$, then substituting $x=2 \al$ and $\eps=2 \de'$ into the same formula, we get the following: 
%
\begin{align*} 
H_1 &\defeq h(1-\al-\de) - h(1-\al) 
= \de - \al\de - \frac{\de^2}{2}  
- \underbrace{\sum_{k=3}^\infty 
\frac{(\al+\de)^k-\al^k}{k(k-1)} }_{S_1 \defeq} ;\\
H_2 &\defeq h(1-2\al-2\de') - h(1-2\al) 
= 2 \de' - 4 \al \de' - 2 (\de')^2 
- \underbrace{\sum_{k=3}^\infty 
2^k \frac{(\al+\de')^k-\al^k}{k(k-1)}}_{S_2 \defeq} .
\end{align*}
%
Therefore, using $\de'=\de-\ga$, 
%
\begin{align} 
\begin{split}
\frac{d}{2} H_2 - (d-1)H_1
&= d\de'-2d\al\de' - d (\de')^2 - \frac{d}{2} S_2\\
&- (d-1) \de 
+ \underbrace{(d-1)\al \de + (d-1)\de^2/2 + (d-1) S_1}_
{\leq d\al \de + d\de^2/2 + d S_1} \\
&\leq \big( -d\al + 1 \big) \de + 
d\left( -\ga -(\de')^2 + \frac{\de^2}{2} + 2\al\ga \right) 
+ d\underbrace{ \left( S_1 - \frac{S_2}{2} \right)}_{\leq 0} ,
\label{eq:bnd3}
\end{split}
\end{align}
%
where the last term is negative due to the following observation: after expansion, we have 
%
\[ S_1 = \sum_{k=3}^\infty \frac{(\al+\de)^k-\al^k}{k(k-1)} 
= \sum_{k=3}^\infty \sum_{i=1}^k a_{k,i} \al^{k-i} \de^i 
\text{ and } 
\frac{S_2}{2} = \sum_{k=3}^\infty \sum_{i=1}^k 
2^{k-1} a_{k,i} \al^{k-i} (\de')^i
\]
%
for some positive coefficients $a_{k,i}$. Due to \eqref{eq:deperde}, $\de^i \leq 2^{k-1} (\de')^i$ holds for any $i \leq k$ and $k \geq 3$, and $S_1 \leq S_2/2$ follows.  

Summing up the bounds \eqref{eq:bnd1}, \eqref{eq:bnd2}, \eqref{eq:bnd3} and adding $2 h(\de)$:
%
\[ B_\de 
\leq 2 h(\de) 
+ \big( -d\al + \log(\al) + 2 \big) \de 
+ d\bigg( 
\underbrace{\ga -(\de')^2 + \frac{\de^2}{2} + 2\al\ga}_
{R \defeq} \bigg) .\]
%
To complete the proof, it suffices to show that $R \leq \Coefd \de^2$ for the last term on the right-hand side. First note that 
%
\[ (\de')^2 = (\de-\ga)^2 \geq \de^2 - 2\de \ga .\]
%
Therefore, using \eqref{eq:ga_bnd}, 
%
\[ R \leq (1+2\al+2\de) \ga - \de^2/2 
\leq \left( \frac{1+2\al+2\de}{1-2\al} 
- \frac{1}{2} \right) \de^2 .\]
%
It remains to show that the coefficient of $\de^2$ is less than $\Coefd$. Note that this coefficient is monotone increasing both in $\al$ and in $\de$, so we may set them to their maximal values, that is, $\al=(2 \log d)/d$ and $\de=3/(2d)$, in which case the coefficient becomes  
%
\[ \frac{1+2\al+2\de}{1-2\al} - \frac{1}{2} 
= \frac{d/2+6 \log d + 3}{d-4 \log d} ,\]
which is indeed less than $\Coefd$ provided that $d \geq 318$.
%
\end{proof}

\end{document}
